\section{数论}\label{sec:number-theory}

模幂、扩展欧几里得、一次不定方程、判素与分解、模平方根，以及可组合的 Universal Euclidean Algorithm。

\subsection{接口}

\apitable{
    \api{qpow<T>(a, b, p)}{返回 $a^b \bmod p$；要求 \code{sizeof(T) <= 8} 且 $b \ge 0$}
    \api{exgcd<T>(a, b, x, y)}{返回 $\gcd(a,b)$，并把一组解写入 \code{x}、\code{y}}
    \api{biequation<T>(a, b, c)}{解 $ax+by=c$，返回 $x$ 最小非负的一组 \code{{x,y}}；无解返回 \code{{-1,-1}}}
    \api{is_prime(i64 n)}{Miller–Rabin 判素；$n<2$ 返回 false，64 位内确定性正确}
    \api{pollard_rho(i64 n)}{返回 $n$ 的一个非平凡因子；要求 $n$ 为合数且 $n>3$}
    \api{factorize(i64 n)}{返回升序质因子（含重数）；\code{n} 为 1 时返回空数组}
    \api{legendre(a, p)}{Legendre 符号 $\left(\frac{a}{p}\right)\in\{-1,0,1\}$，要求 $p$ 为奇素数}
    \api{cipolla(n, p)}{返回 $x$ 使 $x^2\equiv n\pmod p$；无解返回 \code{std::nullopt}}
    \api{uniclidean<D,T,Op>(p, q, r, n, R, U, init, mul)}{返回格路径上各段按 \code{mul} 合并的结果}
}

\subsection{约定}

\code{qpow} 的指数必须非负，负数会在移位循环中死循环，模数不要求素数。\code{biequation} 无解时返回 \code{{-1, -1}}，有解时返回 $x$ 最小非负的一组解。\code{factorize} 只断言 $n>0$：$n=1$ 返回空数组，结果升序且含重数；\code{pollard_rho} 要求输入为合数，内部随机数种子固定，结果可复现。\code{legendre} 与 \code{cipolla} 要求 $p$ 为奇素数，\code{cipolla} 返回其中一根，需要较小根时自行取 $\min(x,p-x)$。

\fullsource{number_theory.hpp}
